\chapter{Rekonesans infrastruktury IT i sieciowej}\label{ch:infrastruktura}

\section{Przestrzeń adresowa IP i systemy autonomiczne}\label{sec:asn-ip}
\subsection{Identyfikacja numerów ASN i bloków adresowych BGP}\label{subsec:asn-bgp}
\subsection{Geolokalizacja i routing sieci szkieletowej}\label{subsec:routing}

\section{Domeny i subdomeny powiązane z jednostką}\label{sec:domeny-subdomeny}
\subsection{Struktura domeny głównej i stref wydziałowych}\label{subsec:strefy-domenowe}
\subsection{Enumeracja subdomen jednostek badawczych i projektowych}\label{subsec:enumeracja-subdomen}

\section{Analiza konfiguracji DNS i bezpieczeństwa poczty}\label{sec:dns}
\subsection{Rekordy zasobów (A, AAAA, CNAME, MX, NS, TXT, SRV)}\label{subsec:rekordy-dns}
\subsection{Mechanizmy ochrony poczty elektronicznej (SPF, DKIM, DMARC)}\label{subsec:ochrona-poczty}

\section{Analiza certyfikatów kryptograficznych TLS}\label{sec:certyfikaty-tls}
\subsection{Eksploracja logów Certificate Transparency (crt.sh)}\label{subsec:ct-logs}
\subsection{Urzędy certyfikacji, parametry kluczy i wygasłe certyfikaty}\label{subsec:analiza-certyfikatow}

\section{Ekspozycja serwerów i usług sieciowych}\label{sec:uslugi-sieciowe}
\subsection{Dane z wyszukiwarek zasobów internetowych (Shodan, Censys)}\label{subsec:shodan-censys}
\subsection{Identyfikacja technologii webowych, nagłówków HTTP i banerów usług}\label{subsec:technologie-web}
\subsection{Systemy uwierzytelniania, portale logowania, VPN i repozytoria kodu}\label{subsec:portale-logowania}

